% s-340-p340-m003; original PDF 340, 341
\Aphorism{s-340-p340-m003}{391}
Im Allgemeinen sind wir als Synthetiker uns als Analytikern voraus; Thätigkeit und Genuss gehen nur in diesem Normalzustande rein und leicht von Statten. Was dies heissen will, erfährt man am besten, wenn eine Zeit lang die Ausnahme herrscht, dass der lebende Mensch vom erkennenden, sich selbst erkennenden eingeholt oder überholt wird, wenn wir durch Abstraction stärker verbrauchen als wir an Concretem nachfüttern. Zu Zwecken des Lebens zerreissen wir den Zusammenhang der Erscheinungen und formen Dinge, Gestalten, Einzelheiten daraus; die Erkenntniss ist gegenwirkend stets bemüht, uns das Ergriffene aus Händen zu nehmen und in den Schmelztiegel zurückzuthun. So lernen wir z. B., indem wir uns isolirte Thatbestände vorführen und Relationen zwischen ihnen suchen; das zunehmende Wissen zeigt dieser Relationen aber so viele, dass ein ganzes Gebiet überschwemmt, jede Grenze verwischt, jede vorläufige Isolirung endgültig aufgehoben wird. Band und Verbundenes ist nicht mehr zu trennen, der Lehrsatz verschmilzt mit dem Beweis; die Sprache versagt, weil alles, was sie vorbringt, einer Correction, einer näheren Bestimmung, einer Angabe einschränkender Bedingungen, eines Hinweises auf Nebenfälle und Nachbargebiete bedarf. Lehren heisst bei diesem Erkenntniss-Stande eigentlich wider besseres Wissen handeln, heisst sich zurückversetzen in eine halbe und unhaltbare Ansicht der Dinge, heisst Identitäten als Sätze, analytische als synthetische Urtheile, mangelnde Gegengründe als Gründe aussprechen. Das Zerreissen der Einheit, um sie nachher künstlich wieder zusammenzufügen, fiel uns das erste Mal leicht, wo es von uns als Function des Lebens ausgeübt wurde; das zweite Mal, wo es der Darstellung unserer Erkenntniss gilt, ist es im Grunde unnatürlich.
\par

